\section{Boundaries}

Reporting order and operational structure routinely create legitimate change
points and clusters. Mail ballots may be reported before precinct ballots;
central-count batches may be uploaded together; recount or adjudication events
may arrive late; scanners and batches may differ greatly in volume. A method
that ignores those facts will produce impressive-looking false positives.

Consuming V-series audit discrepancies does not change the W-series claim. The
V-series records are evidence-bound audit observations. A W-series clustering
report can ask whether their distribution is unusual, but it cannot promote
that unusual distribution into proof of fraud, a certification result, or a
replacement for a risk-limiting audit. The transcript kinds must remain
separate.

The methods are descriptive, not causal. They identify where a sequence changes
or where discrepancies concentrate. The reason for that pattern may be benign,
procedural, data-quality related, or something that merits further review. The
method does not decide among those explanations. It produces a hash-bound,
caveated investigation queue and stops there.

The leading boundary therefore returns at the end: change-point tests and
discrepancy-clustering tests can prioritize investigation. They cannot prove
fraud, certify correctness, or replace risk-limiting audits.
